\documentclass[11pt,a4paper,bibliography=totoc,listof=totoc,pointlessnumbers,open=any]{scrbook}

\usepackage[left=3.0cm,right=2.5cm,top=2.5cm,bottom=3cm]{geometry}	% Seitenränder
\headsep = 1.0cm
\footskip = 1.0cm
                          
\addtokomafont{disposition}{\rmfamily}
\addtokomafont{chapter}{\rmfamily}
\addtokomafont{section}{\rmfamily}
\addtokomafont{subsection}{\rmfamily}
\addtokomafont{subsubsection}{\rmfamily}

\usepackage{setspace} 						% Zeilenabstand 1,5
\onehalfspacing

\usepackage[utf8]{inputenc}					
\usepackage[english]{babel}					% Silbentrennung
\usepackage{amsmath}						% mathematische Formeln
\usepackage{amsfonts}						% mathematische Schriften
\usepackage{amssymb}						% mathematische Symbole
\usepackage{graphicx}						% Einbinden von Graphiken
\usepackage[version=4]{mhchem}
\usepackage{subfigure}						% Ermöglicht ein Subfigure
\usepackage{pdfpages} 						
\usepackage{acronym}						% Symbol- und Abkürzungsverzeichnis
\usepackage[figuresright]{rotating}
\usepackage{subcaption}

%\usepackage[table]{xcolor}
%\usepackage[T1]{fontenc}
%\usepackage{mathptmx}
%\usepackage{colortbl}
%-------MATLAB-------------------------------------------------

%\usepackage{matlab-prettifier}
%\definecolor{number}{gray}{0.6}
%\definecolor{back}{gray}{0.95}

\usepackage{color}                  %\definecolor
\usepackage{listings}               %Programmcodeumgebung für Matlab Code

\usepackage{units} 							% Schreiben von Zahlen mit Einheit mit kleinerem Abstand
\usepackage{fix-cm}							% Schriften größer als \Huge
\usepackage{multirow} 						% Tabelle - Spalten verbinden
\usepackage{array}
\usepackage{here}							% Gleitobjekte "fest setzen"

\usepackage{float}							% ermöglicht Gleitobjekte
\usepackage[font=footnotesize,labelfont=bf]{caption}						% ermöglicht Umbrüche in Bildunterschriften

\usepackage{chngcntr}						% Nummerierung nach Kapitel für Abbildungen, Tabellen, Formeln
\counterwithin{figure}{chapter}
\counterwithin{table}{chapter}
\counterwithin{equation}{chapter}

\usepackage{longtable}						% lange Tabellen über mehrere Seiten
\usepackage{multirow}
\usepackage{booktabs}

\usepackage[headsepline]{scrlayer-scrpage}
\pagestyle{scrheadings}
\automark[section]{chapter}


\usepackage{longtable}						% lange Tabellen über mehrere Seiten

\usepackage[tracking=true]{microtype}		% Anpassen der Abstände zwischen den Buchstaben für Blocksatz

\graphicspath{{Figures/}} 					% Pfad für Bidler
\def\figurename{Figure } 					% keyword for the captions

\DeclareMicrotypeSet*[tracking]{my} %
{ font = */*/*/sc/* }% 
\SetTracking{ encoding = *, shape = sc }{ 45 }% Hier wird festgelegt, dass alle Passagen in Kapitälchen automatisch leicht gesperrt werden. 

\usepackage[colorlinks=false, pdfborder={0 0 0}]{hyperref} % Verlinkung von Inhalts- und Abbildungsverzeichnis 
\usepackage[all]{hypcap} 

\usepackage[numbers]{natbib}				% Paket zur Einbindung des Literaturverzeichnisses 

\usepackage{scrhack}						% Behebt Probleme zwischen den einzelnen Paketen

\begin{document}
\pagestyle{scrheadings}

\newpage
\thispagestyle{empty}
\begin{flushleft}
	\includegraphics[width=0.5\textwidth]{Titel/OTH_Logo_RAI.jpg}
\end{flushleft}

\begin{centering}
	\bigskip 
	\bigskip
	\bigskip
	\bigskip		
	\bigskip
	\huge\textrm{\textbf{A Wiki or overview of international regulatory relations for medical devices }}\\
	\bigskip
	\bigskip
	\bigskip 
	\bigskip 
	\bigskip
		\bigskip
		\bigskip 
		\bigskip
		\bigskip
	\LARGE\textrm{\textbf{Seminar paper}}\\
	\bigskip
	\small\textrm{for the module RFM}\\
	\smallskip
	%\large\textrm{\textbf{\glqq Master of Science (M.Sc.)\grqq}}\\
	%\smallskip
	\small\textrm{at the}\\
	\smallskip
	\large\textrm{\textbf{Ostbayerischen Technischen Hochschule Regensburg}}\\
	\smallskip
	\small\textrm{Fakultät Maschinenbau}\\
	\smallskip
	\small\textrm{Regulatory Affairs and Innovation for Medical Devices Lab}\\
	\smallskip
	\small\textrm{for the degree program}\\
	\smallskip
	\large\textrm{\textbf{Master Medizintechnik}}
	\bigskip 
	\bigskip 
	\bigskip
	\bigskip
	\bigskip 
	\bigskip
	\bigskip 
	\bigskip
	\bigskip 
    \begin{table}[H]
    	\flushleft
    	\renewcommand{\arraystretch}{1.2}
    	\begin{tabular}{ll}
    	\textbf{Authors:} & \textrm{Nova Sassin} \\ 
    	 & \textrm{Matr.-Nr: } 3215867 \\
    
    	\textbf{Supervisor:} & \textrm{Prof. Dr. rer. pol. Max D. Singh }
    	\end{tabular}
    \end{table}
\end{centering}

\newpage
\tableofcontents 

\newpage
\chapter{Responsibilities and authorities in regulatory affairs and their interaction on national, european and U.S. american level}
\section{Authorities and responsibilities in the USA }
\subsection{Legislative Process}

The process of creating and approving laws in the United States involves several steps and the collaboration of various institutions, including Congress, the Senate, and the House of Representatives. As it is a quite complicated procedure, the following section provides a quick overview or summary of the legislative process.


\end{document}
